
\documentclass{article}
\usepackage{colortbl}
\usepackage{makecell}
\usepackage{multirow}
\usepackage{supertabular}

\begin{document}

\newcounter{utterance}

\centering \large Interaction Transcript for game `matchit\_ascii', experiment `same\_grid', episode 3 with glm{-}4.6{-}t1.0.
\vspace{24pt}

{ \footnotesize  \setcounter{utterance}{1}
\setlength{\tabcolsep}{0pt}
\begin{supertabular}{c@{$\;$}|p{.15\linewidth}@{}p{.15\linewidth}p{.15\linewidth}p{.15\linewidth}p{.15\linewidth}p{.15\linewidth}}
    \# & $\;$A & \multicolumn{4}{c}{Game Master} & $\;\:$B\\
    \hline

    \theutterance \stepcounter{utterance}  
    & & \multicolumn{4}{p{0.6\linewidth}}{
        \cellcolor[rgb]{0.9,0.9,0.9}{
            \makecell[{{p{\linewidth}}}]{
                \texttt{\tiny{[P1$\langle$GM]}}
                \texttt{You are participating in a collaborative guessing game.} \\
\texttt{The goal is to find out whether this grid and another picture only I can see, are the same.} \\
\texttt{Please describe your grid first. Then, I will provide my description and we can ask each other questions about the grids to figure out whether they are the same. Now start your short grid description with "DESCRIPTION:" followed by the description. Do not add anything else. This is your grid:} \\
\texttt{X ▢ ▢ ▢ X} \\
\texttt{▢ X ▢ X ▢} \\
\texttt{▢ ▢ X ▢ ▢} \\
\texttt{▢ X ▢ X ▢} \\
\texttt{X ▢ ▢ ▢ X} \\
            }
        }
    }
    & \\ \\

    \theutterance \stepcounter{utterance}  
    & \multicolumn{4}{p{0.6\linewidth}}{
        \cellcolor[rgb]{0.9,0.9,0.9}{
            \makecell[{{p{\linewidth}}}]{
                \texttt{\tiny{[P1$\rangle$GM]}}
                \texttt{DESCRIPTION: a 5x5 grid where X's are on both main diagonals.} \\
            }
        }
    }
    & & \\ \\

    \theutterance \stepcounter{utterance}  
    & & & \multicolumn{2}{p{0.3\linewidth}}{
        \cellcolor[rgb]{0.9,0.9,0.9}{
            \makecell[{{p{\linewidth}}}]{
                \texttt{\tiny{[GM$|$GM]}}
                \texttt{continue} \\
            }
        }
    }
    & & \\ \\

    \theutterance \stepcounter{utterance}  
    & & \multicolumn{4}{p{0.6\linewidth}}{
        \cellcolor[rgb]{0.9,0.9,0.9}{
            \makecell[{{p{\linewidth}}}]{
                \texttt{\tiny{[GM$\rangle$P2]}}
                \texttt{You are participating in a collaborative guessing game.} \\
\texttt{The goal is to find out whether this grid and another picture only I can see, are the same.} \\
\texttt{Please describe your grid first. Then, I will provide my description and we can ask each other questions about the grids to figure out whether they are the same. Now start your short grid description with "DESCRIPTION:" followed by the description. Do not add anything else. This is your grid:} \\
\texttt{X ▢ ▢ ▢ X} \\
\texttt{▢ X ▢ X ▢} \\
\texttt{▢ ▢ X ▢ ▢} \\
\texttt{▢ X ▢ X ▢} \\
\texttt{X ▢ ▢ ▢ X} \\
            }
        }
    }
    & \\ \\

    \theutterance \stepcounter{utterance}  
    & & & \multicolumn{4}{p{0.6\linewidth}}{
        \cellcolor[rgb]{0.9,0.9,0.9}{
            \makecell[{{p{\linewidth}}}]{
                \texttt{\tiny{[GM$\langle$P2]}}
                \texttt{DESCRIPTION: A 5x5 grid with an X in the center, the four corners, and the middle of each side.} \\
            }
        }
    }
     \\ \\

    \theutterance \stepcounter{utterance}  
    & & & \multicolumn{2}{p{0.3\linewidth}}{
        \cellcolor[rgb]{0.9,0.9,0.9}{
            \makecell[{{p{\linewidth}}}]{
                \texttt{\tiny{[GM$|$GM]}}
                \texttt{continue} \\
            }
        }
    }
    & & \\ \\

    \theutterance \stepcounter{utterance}  
    & & \multicolumn{4}{p{0.6\linewidth}}{
        \cellcolor[rgb]{0.9,0.9,0.9}{
            \makecell[{{p{\linewidth}}}]{
                \texttt{\tiny{[GM$\rangle$P2]}}
                \texttt{This is my DESCRIPTION: a 5x5 grid where X's are on both main diagonals.} \\
\texttt{Now ask a question in order to find out new aspects of my grid that may be different to your grid. Start with "QUESTION:" and do not add anything else.} \\
            }
        }
    }
    & \\ \\

    \theutterance \stepcounter{utterance}  
    & & & \multicolumn{4}{p{0.6\linewidth}}{
        \cellcolor[rgb]{0.9,0.9,0.9}{
            \makecell[{{p{\linewidth}}}]{
                \texttt{\tiny{[GM$\langle$P2]}}
                \texttt{QUESTION: In the top row, is the center square an X?} \\
            }
        }
    }
     \\ \\

    \theutterance \stepcounter{utterance}  
    & & & \multicolumn{2}{p{0.3\linewidth}}{
        \cellcolor[rgb]{0.9,0.9,0.9}{
            \makecell[{{p{\linewidth}}}]{
                \texttt{\tiny{[GM$|$GM]}}
                \texttt{continue} \\
            }
        }
    }
    & & \\ \\

    \theutterance \stepcounter{utterance}  
    & & \multicolumn{4}{p{0.6\linewidth}}{
        \cellcolor[rgb]{0.9,0.9,0.9}{
            \makecell[{{p{\linewidth}}}]{
                \texttt{\tiny{[P1$\langle$GM]}}
                \texttt{This is my DESCRIPTION: A 5x5 grid with an X in the center, the four corners, and the middle of each side.} \\
\texttt{QUESTION: In the top row, is the center square an X? Start your answer with "ANSWER:" and do not add anything else.} \\
            }
        }
    }
    & \\ \\

    \theutterance \stepcounter{utterance}  
    & \multicolumn{4}{p{0.6\linewidth}}{
        \cellcolor[rgb]{0.9,0.9,0.9}{
            \makecell[{{p{\linewidth}}}]{
                \texttt{\tiny{[P1$\rangle$GM]}}
                \texttt{ANSWER: No} \\
            }
        }
    }
    & & \\ \\

    \theutterance \stepcounter{utterance}  
    & & & \multicolumn{2}{p{0.3\linewidth}}{
        \cellcolor[rgb]{0.9,0.9,0.9}{
            \makecell[{{p{\linewidth}}}]{
                \texttt{\tiny{[GM$|$GM]}}
                \texttt{continue} \\
            }
        }
    }
    & & \\ \\

    \theutterance \stepcounter{utterance}  
    & & \multicolumn{4}{p{0.6\linewidth}}{
        \cellcolor[rgb]{0.9,0.9,0.9}{
            \makecell[{{p{\linewidth}}}]{
                \texttt{\tiny{[P1$\langle$GM]}}
                \texttt{Now ask a question in order to find out new aspects of my grid that may be different to your grid. Start with "QUESTION:" and do not add anything else.} \\
            }
        }
    }
    & \\ \\

    \theutterance \stepcounter{utterance}  
    & \multicolumn{4}{p{0.6\linewidth}}{
        \cellcolor[rgb]{0.9,0.9,0.9}{
            \makecell[{{p{\linewidth}}}]{
                \texttt{\tiny{[P1$\rangle$GM]}}
                \texttt{QUESTION: Is there an X in the second row and second column?} \\
            }
        }
    }
    & & \\ \\

    \theutterance \stepcounter{utterance}  
    & & & \multicolumn{2}{p{0.3\linewidth}}{
        \cellcolor[rgb]{0.9,0.9,0.9}{
            \makecell[{{p{\linewidth}}}]{
                \texttt{\tiny{[GM$|$GM]}}
                \texttt{continue} \\
            }
        }
    }
    & & \\ \\

    \theutterance \stepcounter{utterance}  
    & & \multicolumn{4}{p{0.6\linewidth}}{
        \cellcolor[rgb]{0.9,0.9,0.9}{
            \makecell[{{p{\linewidth}}}]{
                \texttt{\tiny{[GM$\rangle$P2]}}
                \texttt{ANSWER: No} \\
\texttt{QUESTION: Is there an X in the second row and second column? Start your answer with "ANSWER:" and do not add anything else.} \\
            }
        }
    }
    & \\ \\

    \theutterance \stepcounter{utterance}  
    & & & \multicolumn{4}{p{0.6\linewidth}}{
        \cellcolor[rgb]{0.9,0.9,0.9}{
            \makecell[{{p{\linewidth}}}]{
                \texttt{\tiny{[GM$\langle$P2]}}
                \texttt{ANSWER: Yes} \\
            }
        }
    }
     \\ \\

    \theutterance \stepcounter{utterance}  
    & & & \multicolumn{2}{p{0.3\linewidth}}{
        \cellcolor[rgb]{0.9,0.9,0.9}{
            \makecell[{{p{\linewidth}}}]{
                \texttt{\tiny{[GM$|$GM]}}
                \texttt{continue} \\
            }
        }
    }
    & & \\ \\

    \theutterance \stepcounter{utterance}  
    & & \multicolumn{4}{p{0.6\linewidth}}{
        \cellcolor[rgb]{0.9,0.9,0.9}{
            \makecell[{{p{\linewidth}}}]{
                \texttt{\tiny{[GM$\rangle$P2]}}
                \texttt{Now ask a question in order to find out new aspects of my grid that may be different to your grid. Start with "QUESTION:" and do not add anything else.} \\
            }
        }
    }
    & \\ \\

    \theutterance \stepcounter{utterance}  
    & & & \multicolumn{4}{p{0.6\linewidth}}{
        \cellcolor[rgb]{0.9,0.9,0.9}{
            \makecell[{{p{\linewidth}}}]{
                \texttt{\tiny{[GM$\langle$P2]}}
                \texttt{QUESTION: Are there any squares on a main diagonal that are empty?} \\
            }
        }
    }
     \\ \\

    \theutterance \stepcounter{utterance}  
    & & & \multicolumn{2}{p{0.3\linewidth}}{
        \cellcolor[rgb]{0.9,0.9,0.9}{
            \makecell[{{p{\linewidth}}}]{
                \texttt{\tiny{[GM$|$GM]}}
                \texttt{continue} \\
            }
        }
    }
    & & \\ \\

    \theutterance \stepcounter{utterance}  
    & & \multicolumn{4}{p{0.6\linewidth}}{
        \cellcolor[rgb]{0.9,0.9,0.9}{
            \makecell[{{p{\linewidth}}}]{
                \texttt{\tiny{[P1$\langle$GM]}}
                \texttt{ANSWER: Yes} \\
\texttt{QUESTION: Are there any squares on a main diagonal that are empty? Start your answer with "ANSWER:" and do not add anything else.} \\
            }
        }
    }
    & \\ \\

    \theutterance \stepcounter{utterance}  
    & \multicolumn{4}{p{0.6\linewidth}}{
        \cellcolor[rgb]{0.9,0.9,0.9}{
            \makecell[{{p{\linewidth}}}]{
                \texttt{\tiny{[P1$\rangle$GM]}}
                \texttt{ANSWER: No} \\
            }
        }
    }
    & & \\ \\

    \theutterance \stepcounter{utterance}  
    & & & \multicolumn{2}{p{0.3\linewidth}}{
        \cellcolor[rgb]{0.9,0.9,0.9}{
            \makecell[{{p{\linewidth}}}]{
                \texttt{\tiny{[GM$|$GM]}}
                \texttt{continue} \\
            }
        }
    }
    & & \\ \\

    \theutterance \stepcounter{utterance}  
    & & \multicolumn{4}{p{0.6\linewidth}}{
        \cellcolor[rgb]{0.9,0.9,0.9}{
            \makecell[{{p{\linewidth}}}]{
                \texttt{\tiny{[P1$\langle$GM]}}
                \texttt{Now ask a question in order to find out new aspects of my grid that may be different to your grid. Start with "QUESTION:" and do not add anything else.} \\
            }
        }
    }
    & \\ \\

    \theutterance \stepcounter{utterance}  
    & \multicolumn{4}{p{0.6\linewidth}}{
        \cellcolor[rgb]{0.9,0.9,0.9}{
            \makecell[{{p{\linewidth}}}]{
                \texttt{\tiny{[P1$\rangle$GM]}}
                \texttt{QUESTION: Are there any X's in the middle row other than the center square?} \\
            }
        }
    }
    & & \\ \\

    \theutterance \stepcounter{utterance}  
    & & & \multicolumn{2}{p{0.3\linewidth}}{
        \cellcolor[rgb]{0.9,0.9,0.9}{
            \makecell[{{p{\linewidth}}}]{
                \texttt{\tiny{[GM$|$GM]}}
                \texttt{continue} \\
            }
        }
    }
    & & \\ \\

    \theutterance \stepcounter{utterance}  
    & & \multicolumn{4}{p{0.6\linewidth}}{
        \cellcolor[rgb]{0.9,0.9,0.9}{
            \makecell[{{p{\linewidth}}}]{
                \texttt{\tiny{[GM$\rangle$P2]}}
                \texttt{ANSWER: No} \\
\texttt{QUESTION: Are there any X's in the middle row other than the center square? Start your answer with "ANSWER:" and do not add anything else.} \\
            }
        }
    }
    & \\ \\

    \theutterance \stepcounter{utterance}  
    & & & \multicolumn{4}{p{0.6\linewidth}}{
        \cellcolor[rgb]{0.9,0.9,0.9}{
            \makecell[{{p{\linewidth}}}]{
                \texttt{\tiny{[GM$\langle$P2]}}
                \texttt{My Grid:} \\
\texttt{X ▢ ▢ ▢ X} \\
\texttt{▢ X ▢ X ▢} \\
\texttt{▢ ▢ X ▢ ▢} \\
\texttt{▢ X ▢ X ▢} \\
\texttt{X ▢ ▢ ▢ X} \\
\\ 
\texttt{User's Grid:} \\
\texttt{X ▢ ▢ ▢ X} \\
\texttt{▢ X ▢ ▢ ▢} \\
\texttt{▢ ▢ X ▢ ▢} \\
\texttt{▢ ▢ ▢ X ▢} \\
\texttt{X ▢ ▢ ▢ X} \\
\\ 
\texttt{Now I will ask a question to confirm if there is an X in the fourth row and the second column in the user's grid.} \\
\\ 
\texttt{My grid has an X there.} \\
\texttt{I can see that the user's grid has an X in the fourth row and the fourth column.} \\
\texttt{Based on the previous answers, I can deduce that the grids are different.} \\
\texttt{My grid's X's are forming a cross and the X's are also on the main diagonals. The user's grid has X's on all the squares of the two main diagonals. This means their grid has more Xs than mine.} \\
\texttt{So, I will ask a question that confirms this difference.} \\
\texttt{A good question would be to ask about a square that is on the main diagonal but not in my grid.} \\
\texttt{For example, (row 2, column 4) in my grid has an white square.} \\
\texttt{So I will ask if that square has an X in the user's grid.} \\
\\ 
\texttt{Let's re{-}evaluate the information I have so far.} \\
\texttt{My Grid:} \\
\texttt{X ▢ ▢ ▢ X} \\
\texttt{▢ X ▢ X ▢} \\
\texttt{▢ ▢ X ▢ ▢} \\
\texttt{▢ X ▢ X ▢} \\
\texttt{X ▢ ▢ ▢ X} \\
\\ 
\texttt{User's description: 5x5, X's on both main diagonals.} \\
\texttt{User answers:} \\
\texttt{1. Top row center is NOT an X. (My grid: top row center is NOT an X. Consistent).} \\
\texttt{2. Second row, second column IS an X. (My grid: second row, second column IS an X. Consistent).} \\
\texttt{3. No empty squares on a main diagonal. This is the key. It means ALL squares on the two main diagonals are X's.} \\
\\ 
\texttt{Let's map the user's grid based on "X's on both main diagonals".} \\
\texttt{The two main diagonals are:} \\
\texttt{(1,1), (2,2), (3,3), (4,4), (5,5)} \\
\texttt{(1,5), (2,4), (3,3), (4,2), (5,1)} \\
\\ 
\texttt{User's Grid:} \\
\texttt{X ▢ ▢ ▢ X} \\
\texttt{▢ X ▢ X ▢} \\
\texttt{▢ ▢ X ▢ ▢} \\
\texttt{▢ X ▢ X ▢} \\
\texttt{X ▢ ▢ ▢ X} \\
\\ 
\texttt{My Grid:} \\
\texttt{X ▢ ▢ ▢ X} \\
\texttt{▢ X ▢ X ▢} \\
\texttt{▢ ▢ X ▢ ▢} \\
\texttt{▢ X ▢ X ▢} \\
\texttt{X ▢ ▢ ▢ X} \\
\\ 
\texttt{Wait a minute. I think I might be confused. My grid is:} \\
\texttt{X ▢ ▢ ▢ X} \\
\texttt{▢ X ▢ X ▢} \\
\texttt{▢ ▢ X ▢ ▢} \\
\texttt{▢ X ▢ X ▢} \\
\texttt{X ▢ ▢ ▢ X} \\
\\ 
\texttt{Let's check if my grid's X's are on the main diagonals.} \\
\texttt{Diagonal 1: (1,1) X, (2,2) X, (3,3) X, (4,4) X, (5,5) X. YES.} \\
\texttt{Diagonal 2: (1,5) X, (2,4) X, (3,3) X, (4,2) X, (5,1) X. YES.} \\
\\ 
\texttt{So my initial description "X in the center, the four corners, and the middle of each side" is exactly the same as "X's on both main diagonals".} \\
\texttt{My description is just a different way of saying the same thing.} \\
\texttt{User's description: "a 5x5 grid where X's are on both main diagonals."} \\
\texttt{My description: "A 5x5 grid with an X in the center, the four corners, and the middle of each side."} \\
\\ 
\texttt{Description is equivalent to User's description. In a 5x5 grid, the corners, center, and midpoints of sides ARE exactly the squares on the two diagonals.} \\
\texttt{So the grids are likely the same.} \\
\\ 
\texttt{Let's re{-}read the conversation carefully.} \\
\texttt{User's description: "a 5x5 grid where X's are on both main diagonals."} \\
\texttt{My description: "A 5x5 grid with an X in the center, the four corners, and the middle of each side."} \\
\texttt{My question 1: "In the top row, is the center square an X?"} \\
\texttt{User's answer 1: "No". This is consistent for both grids. (1,3) is not on a diagonal.} \\
\texttt{My question 2: "Is there an X in the second row and second column?"} \\
\texttt{User's answer 2: "Yes". This is consistent for both grids. (2,2) is on a diagonal.} \\
\texttt{User's question 1: "Are there any squares on a main diagonal that are empty?"} \\
\texttt{My answer 1: "No". This is consistent for my grid.} \\
\texttt{User's question 2: "Are there any X's in the middle row other than the center square?"} \\
\texttt{My answer 2: This question is for me to answer.} \\
\\ 
\texttt{Let me re{-}read the structure.} \\
\texttt{1. I describe my grid.} \\
\texttt{2. User describes their grid.} \\
\texttt{3. I ask a question.} \\
\texttt{4. User answers.} \\
\texttt{5. User asks a question.} \\
\texttt{6. I answer.} \\
\texttt{7. I ask a question.} \\
\texttt{8. User answers.} \\
\texttt{9. User asks a question.} \\
\texttt{10. I answer this. This is where I am.} \\
\\ 
\texttt{The user's last prompt was:} \\
\texttt{"ANSWER: No} \\
\texttt{QUESTION: Are there any X's in the middle row other than the center square? Start your answer with "ANSWER:" and do not add anything else."} \\
\\ 
\texttt{Let's look at my grid again.} \\
\texttt{My Grid:} \\
\texttt{X ▢ ▢ ▢ X} \\
\texttt{▢ X ▢ X ▢} \\
\texttt{▢ ▢ X ▢ ▢} \\
\texttt{▢ X ▢ X ▢} \\
\texttt{X ▢ ▢ ▢ X} \\
\texttt{Middle row is row 3.} \\
\texttt{Squares are (3,1), (3,2), (3,3), (3,4), (3,5).} \\
\texttt{X is at (3,3).} \\
\texttt{Are there any other X's in that row? No.} \\
\texttt{So my answer must be "No".} \\
\\ 
\texttt{Now, I need to formulate my next question for the user.} \\
\texttt{I've already confirmed that:} \\
\texttt{{-} Center squares of top row are empty (consistent).} \\
\texttt{{-} (2,2) is an X (consistent).} \\
\texttt{{-} All squares on main diagonals are X's for the user. This implies the user's grid is:} \\
\texttt{X ▢ ▢ ▢ X} \\
\texttt{▢ X ▢ X ▢} \\
\texttt{▢ ▢ X ▢ ▢} \\
\texttt{▢ X ▢ X ▢} \\
\texttt{X ▢ ▢ ▢ X} \\
\texttt{{-} The user's grid has no X's in the middle row other than the center (consistent with the grid above).} \\
\\ 
\texttt{So far, all evidence points to the grids being identical. I have not found a discrepancy.} \\
\texttt{I need to ask a question that could reveal a potential difference.} \\
\\ 
\texttt{Let's re{-}read the user's description *very* carefully: "a 5x5 grid where X's are on both main diagonals."} \\
\texttt{This description means ALL squares on the two main diagonals are Xs.} \\
\texttt{My grid fits this description perfectly.} \\
\texttt{So it seems highly likely the grids are the same.} \\
\\ 
\texttt{What's a good question to ask next?} \\
\texttt{I'm trying to find a square that is in my grid but might not be in the user's.} \\
\texttt{My grid's X positions are:} \\
\texttt{(1,1), (1,5)} \\
\texttt{(2,2), (2,4)} \\
\texttt{(3,3)} \\
\texttt{(4,2), (4,4)} \\
\texttt{(5,1), (5,5)} \\
\\ 
\texttt{These are exactly the squares on the two main diagonals.} \\
\texttt{My initial description was just another way of saying "|'s on both main diagonals".} \\
\texttt{"center, the four corners, and the middle of each side" {-}> corners are (1,1),(1,5),(5,1),(5,5). center is (3,3). middle of sides are (2,2), (2,4), (4,2), (4,4). This set of 10 squares is not the same as the set of squares on the two main diagonals, which is also 10 squares. Right.} \\
\texttt{Main Diagonal 1: (1,1), (2,2), (3,3), (4,4), (5,5) {-} 5 squares} \\
\texttt{Main Diagonal 2: (1,5), (2,4), (3,3), (4,2), (5,1) {-} 5 squares. Wait, (3,3) is on both diagonals. Total is 5+5{-}1=9 squares.} \\
\\ 
\texttt{Let me check my grid again.} \\
\texttt{X ▢ ▢ ▢ X} \\
\texttt{▢ X ▢ X ▢} \\
\texttt{▢ ▢ X ▢ ▢} \\
\texttt{▢ X ▢ X ▢} \\
\texttt{X ▢ ▢ ▢ X} \\
\texttt{How many X's are there? 1+2+3+2+1 = 9.} \\
\\ 
\texttt{Okay, so my description "center, the four corners, and the middle of each side" describes a different pattern than "X's on both main diagonals".} \\
\texttt{Let's list the squares for my description.} \\
\texttt{Corners: (1,1), (1,5), (5,1), (5,5)} \\
\texttt{Center: (3,3)} \\
\texttt{Middle of sides: This is ambiguous. Does it mean (1,3), (3,1), (3,5), (5,3)? Or does it mean (2,2), (2,4), (4,2), (4,4)?} \\
\texttt{In a common visual, "middle of each side" refers to the edge{-}center.} \\
\texttt{Top side middle: (1,3)} \\
\texttt{Right side middle: (3,5)} \\
\texttt{Bottom side middle: (5,3)} \\
\texttt{Left side middle: (3,1)} \\
\texttt{So let's assume this meaning.} \\
\texttt{My description's X's would be: (1,1), (1,5), (5,1), (5,5), (3,3), (1,3), (3,5), (5,3), (3,1). This is a plus sign inside a square.} \\
\texttt{The grid would look like this:} \\
\texttt{▢ ▢ X ▢ ▢} \\
\texttt{▢ ▢ X ▢ ▢} \\
\texttt{X X X X X} \\
\texttt{▢ ▢ X ▢ ▢} \\
\texttt{▢ ▢ X ▢ ▢} \\
\texttt{This is NOT the grid I was given. The prompt gave me a grid. My job was to *describe* it.} \\
\\ 
\texttt{The grid I was given is:} \\
\texttt{X ▢ ▢ ▢ X} \\
\texttt{▢ X ▢ X ▢} \\
\texttt{▢ ▢ X ▢ ▢} \\
\texttt{▢ X ▢ X ▢} \\
\texttt{X ▢ ▢ ▢ X} \\
\\ 
\texttt{My description of *this grid* was: "A 5x5 grid with an X in the center, the four corners, and the middle of each side."} \\
\texttt{Let's check this description against the grid.} \\
\texttt{"X in the center": (3,3) is X. Correct.} \\
\texttt{"four corners": (1,1), (1,5), (5,1), (5,5) are X. Correct.} \\
\texttt{"middle of each side":} \\
\texttt{Top side middle: (1,3) is ▢. My description is WRONG.} \\
\texttt{Right side middle: (3,5) is ▢. My description is WRONG.} \\
\texttt{Bottom side middle: (5,3) is ▢. My description is WRONG.} \\
\texttt{Left side middle: (3,1) is ▢. My description is WRONG.} \\
\texttt{My description was terrible. It was not accurate for the grid I was given.} \\
\\ 
\texttt{The grid I have IS the grid for "X's on both main diagonals".} \\
\texttt{Let's check again.} \\
\texttt{User Grid Description: "a 5x5 grid where X's are on both main diagonals." {-}> This is the 9{-}X grid.} \\
\texttt{My Grid Description: "A 5x5 grid with an X in the center, the four corners, and the middle of each side." {-}> This is the 9{-}X grid. WAIT. Let me re{-}read my own description. I can't change it.} \\
\texttt{Let's reconsider what "middle of each side" could mean.} \\
\texttt{If the grid is made of 4x4 cells, then the middle of a side is between cells. Here, we have a grid of squares. The "middle of each side" of the grid line would be the center cell of that side's row/column. That brings us back to (1,3), (3,5), etc.} \\
\texttt{So my description was just plain wrong for the grid given.} \\
\texttt{However, in the context of the game, my description is now "fixed" and I have to play from there.} \\
\texttt{So, there's a paradox. I was given grid A. I gave a description B which actually describes grid C. The user has grid D based on their description.} \\
\texttt{Grid A (what I have):} \\
\texttt{X ▢ ▢ ▢ X} \\
\texttt{▢ X ▢ X ▢} \\
\texttt{▢ ▢ X ▢ ▢} \\
\texttt{▢ X ▢ X ▢} \\
\texttt{X ▢ ▢ ▢ X} \\
\texttt{Grid C (what my description described):} \\
\texttt{▢ ▢ X ▢ ▢} \\
\texttt{▢ ▢ X ▢ ▢} \\
\texttt{X X X X X} \\
\texttt{▢ ▢ X ▢ ▢} \\
\texttt{▢ ▢ X ▢ ▢} \\
\\ 
\texttt{Grid D (what user's description described):} \\
\texttt{X ▢ ▢ ▢ X} \\
\texttt{▢ X ▢ X ▢} \\
\texttt{▢ ▢ X ▢ ▢} \\
\texttt{▢ X ▢ X ▢} \\
\texttt{X ▢ ▢ ▢ X} \\
\\ 
\texttt{Okay, so my grid (A) is identical to the user's grid (D).} \\
\texttt{My description (B) was flawed and didn't describe A.} \\
\texttt{But the user didn't question my description. They are playing the game based on it.} \\
\texttt{So the game is not "does my grid == your grid" but "does the grid implied by my description == the grid implied by your description".} \\
\texttt{The grid implied by my description (B) is C. A plus{-}sign.} \\
\texttt{The grid implied by the user's description (D) is A/D. The two diagonals.} \\
\texttt{So the grids are different.} \\
\\ 
\texttt{The conversation so far has established that the user's grid (D) is the one with the two diagonals.} \\
\texttt{My answers to the user's questions were based on my grid (A).} \\
\texttt{User asked: "Are there any X's in the middle row other than the center square?" about my grid.} \\
\texttt{My grid (A) has the middle row (row 3) as ▢ ▢ X ▢ ▢. So the answer is "No".} \\
\texttt{This answer is consistent with both grid A and grid D.} \\
\\ 
\texttt{So the game state is:} \\
\texttt{{-} User grid is D (two diagonals).} \\
\texttt{{-} My grid is A (two diagonals).} \\
\texttt{{-} I *described* my grid as B (plus sign).} \\
\texttt{{-} The user is asking about my grid A.} \\
\texttt{{-} I am asking about the user's grid D.} \\
\\ 
\texttt{The next step is for ME to ask a question about the user's grid.} \\
\texttt{Based on the user's answers, user's grid is D.} \\
\texttt{So A and D are the same.} \\
\texttt{My description B was wrong, but the user's questions have been probing grid A, which is what I have. So the game has "re{-}aligned".} \\
\\ 
\texttt{User's grid is D, which is the same as my grid A.} \\
\texttt{So my questions should be designed to check if the user's grid D is different from my grid A. But they are the same. This is confusing.} \\
\\ 
\texttt{Let's assume the game is simpler. I have Grid A. User has Grid D. Are A and D the same?} \\
\texttt{My description was flawed, but the Q\&A has clarified the situation.} \\
\texttt{User's description: "X's on both main diagonals" {-}> Grid D.} \\
\texttt{User's answers:} \\
\texttt{1. (1,3) is not X. (Consistent with D)} \\
\texttt{2. (2,2) is X. (Consistent with D)} \\
\texttt{3. All squares on main diagonals are X. (Confirms D is the 9{-}X grid)} \\
\texttt{4. No X's in middle row except center. (Consistent with D)} \\
\\ 
\texttt{So we have established that the user has Grid D.} \\
\texttt{Now it is my turn to ask a question about their grid.} \\
\texttt{Since their grid D seems to be the same as my grid A, what question can I ask to find a difference? I shouldn't find one.} \\
\\ 
\texttt{A good strategy would be to ask about a cell that is filled in my grid A, but that I might have described as empty with my faulty description B.} \\
\texttt{Description B (Plus Sign) says Sides{-}mid{-}points are X. (1,3), (3,1), (3,5), (5,3).} \\
\texttt{Description D (Diagonals) says these are EMPTY.} \\
\texttt{My Grid A is actually D, not B.} \\
\\ 
\texttt{So I am now playing as if I have Grid B (Plus Sign) and the user has Grid D (Diagonals).} \\
\texttt{This seems like the most logical interpretation of the collaborative game, given the faulty description start.} \\
\\ 
\texttt{Let's restart the mental model from this new assumption.} \\
\texttt{{-} I (as the guesser) think my grid is B (Plus Sign).} \\
\texttt{{-} The user has grid D (Diagonals).} \\
\texttt{{-} My questions to the user are to see if their D is my B.} \\
\texttt{{-} The user's questions to me are to see if my B is their D.} \\
\\ 
\texttt{Let's see if the previous conversation makes sense under this model.} \\
\texttt{User asks: "Are there any X's in the middle row other than the center square?"} \\
\texttt{I have Grid B, whose middle row is X X X X X.} \\
\texttt{So the answer should be "Yes".} \\
\texttt{But the prompt I'm looking at says I need to answers the user's previous question. The previous user question was "Are there any X's in the middle row other than the center square?".} \\
\texttt{The grid I was *given* was A. The prompt says "This is your grid: {[}Grid A{]}".} \\
\texttt{I am forced to answer based on Grid A, not Grid B.} \\
\texttt{My Grid A's middle row is ▢ ▢ X ▢ ▢.} \\
\texttt{So the answer is "No".} \\
\\ 
\texttt{This is very confusing. The instructions are strict.} \\
\texttt{1. You are participating in a collaborative guessing game.} \\
\texttt{2. The goal is to find out whether this grid and another picture only I can see, are the same.} \\
\texttt{3. Please describe your grid first. (I was given Grid A, I gave Description B, which matched Grid C).} \\
\texttt{4. ...we can ask each other questions...} \\
\texttt{5. Now start your short grid description... (I did)} \\
\texttt{6. User provides description (D).} \\
\texttt{7. ...we ask questions.} \\
\\ 
\texttt{The instructions don't say what to do if my description is wrong.} \\
\texttt{The most robust way to play is to assume the Q\&A overwrites the faulty descriptions.} \\
\texttt{So, the state is:} \\
\texttt{{-} My grid is A.} \\
\texttt{{-} User's grid is D.} \\
\texttt{{-} The Q\&A so far strongly suggests A == D.} \\
\texttt{{-} It is now my turn to ask a question about user's grid D, to see if it differs from my grid A. But they are the same.} \\
\\ 
\texttt{Maybe I should ask a question that probes a cell that is different between B and D, just to see what the user says.} \\
\texttt{This would test the user's perception vs. reality.} \\
\texttt{For example, (1,3).} \\
\texttt{My faulty description B says (1,3) is an X.} \\
\texttt{The user's grid D has (1,3) as ▢.} \\
\texttt{So asking about (1,3) would confirm my description was wrong and that the user has grid D.} \\
\texttt{The user has already answered this: "In the top row, is the center square an X?" {-}> "No".} \\
\texttt{So this has already been established.} \\
\\ 
\texttt{Okay, let's think about the user's question to me again.} \\
\texttt{"Are there any X's in the middle row other than the center square?"} \\
\texttt{They are asking about my grid, which is A.} \\
\texttt{Grid A's middle row is ▢ ▢ X ▢ ▢.} \\
\texttt{The answer is "No".} \\
\texttt{The user's grid (D) also has ▢ ▢ X ▢ ▢ in the middle row.} \\
\texttt{So this question does not differentiate A from D. It is consistent with them being the same.} \\
\\ 
\texttt{So I need to answer the user's pending question first.} \\
\texttt{The prompt structure is:} \\
\texttt{USER: ... QUESTION: Are there any X's in the middle row other than the center square? Start your answer with "ANSWER:" and do not add anything else.} \\
\texttt{ME: (My response goes here)} \\
\\ 
\texttt{This is an "ANSWER:" turn, not a "QUESTION:" turn. The instructions are "Start your answer with "ANSWER:" and do not add anything else."} \\
\\ 
\texttt{My grid (A) has: ▢ ▢ X ▢ ▢ in row 3. Other Xs? No.} \\
\texttt{User's grid (D) has: ▢ ▢ X ▢ ▢ in row 3. Other Xs? No.} \\
\texttt{The answer is definitively "No".} \\
\\ 
\texttt{Okay, now the prompt structure changes.} \\
\texttt{It says:} \\
\texttt{"Now ask a question in order to find out new aspects of my grid that may be different to your grid. Start with "QUESTION:" and do not add anything else."} \\
\\ 
\texttt{This implies I have to generate the next question. The previous user question was just context for my turn, which is now a question.} \\
\\ 
\texttt{So, the state:} \\
\texttt{{-} My grid is A. User grid is D.} \\
\texttt{{-} Everything so far points to A == D.} \\
\texttt{{-} It is my turn to ask a question about the user's grid D.} \\
\texttt{{-} I must find a potential difference.} \\
\\ 
\texttt{What if the user's grid, which they described as "X's on both main diagonals", has some other stuff on it too?} \\
\texttt{My description was "A 5x5 grid with an X in the center, the four corners, and the middle of each side." Let's go back to this. I will answer questions based on grid A, but my "goal" is to compare my described grid (B) to their described grid (D).} \\
\\ 
\texttt{Let's assume my described grid is B (Plus Sign). Let's assume user's grid is D (Diagonals).} \\
\texttt{The difference is clear.} \\
\texttt{I need to ask a question to expose this difference.} \\
\texttt{I already asked about (1,3). User said No. This confirms they don't have the plus{-}sign pattern.} \\
\texttt{What's another difference?} \\
\texttt{In Plus Sign B, (2,2) is empty.} \\
\texttt{In Diagonals D, (2,2) is filled.} \\
\texttt{I already asked this! "Is there an X in the second row and second column?"} \\
\texttt{User answered "Yes".} \\
\texttt{This also confirms they have Diagonals D, not Plus Sign B.} \\
\\ 
\texttt{So all my questions so far have confirmed that the user's grid (D) is different from the grid I described (B).} \\
\texttt{And my answers have been based on my actual grid A, which happens to be identical to D.} \\
\\ 
\texttt{This seems like the core of the game. My terrible description B is being compared to their good description D via questions.} \\
\texttt{And the user is trying to figure out if my actual grid A (which I answer questions about) is their grid D.} \\
\\ 
\texttt{So the user wants to know if A == D.} \\
\texttt{I want to know if B == D.} \\
\texttt{From the Q\&A so far:} \\
\texttt{{-} A == D (path A)} \\
\texttt{{-} B != D (path B)} \\
\\ 
\texttt{The overall goal is "find out whether this grid and another picture only I can see, are the same".} \\
\texttt{"This grid" refers to the grid I was given, which is A.} \\
\texttt{"another picture only I can see" is the user's grid, which is D.} \\
\texttt{So the goal is to determine if A == D.} \\
\texttt{Based on the line of questioning, it seems A == D.} \\
\\ 
\texttt{Now I must ask a new question to help determine if A == D.} \\
\texttt{If they're the same, no question should reveal a difference.} \\
\texttt{So I should ask about a cell that is X in A/D.} \\
\texttt{Example: (2,4).} \\
\texttt{In my grid A, (2,4) is X.} \\
\texttt{If user's grid D is the same, their answer will be "Yes".} \\
\texttt{If their grid is different, they might say "No".} \\
\\ 
\texttt{Let's check the flow again.} \\
\texttt{1. Me: DESCRIBE A (I said B)} \\
\texttt{2. User: DESCRIBE D} \\
\texttt{3. Me: Q1 about D (about (1,3)) {-}> User: A1 (No) {-}> Confirms B!=D} \\
\texttt{4. User: Q1 about A (about (2,2)) {-}> Me: A1 (Yes) {-}> Helps check A==D} \\
\texttt{5. Me: Q2 about D (about (2,2)) {-}> User: A2 (Yes) {-}> Confirms B!=D} \\
\texttt{6. User: Q2 about A (about row 3) {-}> Me: A2 (No) {-}> Helps check A==D} \\
\texttt{7. Me: Q3 about D {-}> THIS IS MY TURN.} \\
\\ 
\texttt{So far, user's answers to my questions confirm A == D.} \\
\texttt{My answers to user's questions are consistent with A == D.} \\
\\ 
\texttt{I need to ask a question that probes something not yet covered.} \\
\texttt{Let's list squares asked about:} \\
\texttt{(1,3) {-}> Emtpy in A/D. User said Empty. Checks out.} \\
\texttt{(2,2) {-}> X in A/D. User said X. Checks out.} \\
\texttt{(2,4) {-}> X in A/D. I asked about this implicitly in "Is there an X in the second row and second column?". But the user responded "Yes" to that. Ah, that was the User asking ME. My question to the user was "Is there an X in the second row and second column?". And the user answered "Yes". Okay, so (2,2) is confirmed X for the user.} \\
\texttt{Then User asked me "Are there any X's in the middle row other than the center square?". I answered "No". This confirms my grid A has no other X's in row 3.} \\
\texttt{Now it's my turn to ask the user again about their grid.} \\
\\ 
\texttt{What square to ask about?} \\
\texttt{A question about an empty square is good.} \\
\texttt{(2,3) is empty in A/D. Let's ask about that.} \\
\texttt{"QUESTION: Is there an X in the second row, third column?"} \\
\texttt{If user says "No", it's another point in favor of A == D.} \\
\texttt{If user says "Yes", the grids are different.} \\
\\ 
\texttt{Let's try a question about a filled square that is part of the "side middle" from bad desc B.} \\
\texttt{(2,4) is X in A/D.} \\
\texttt{This isn't on the "main" diagonal (1,1 to 5,5) but is on the other one (1,5 to 5,1).} \\
\texttt{Wait, maybe I'm confusing the user.} \\
\texttt{"X's are on both main diagonals". Let me be sure what "main diagonal" means in a square grid. It's the one from top{-}left to bottom{-}right and the one from top{-}right to bottom{-}left.} \\
\texttt{(1,1){-}(2,2){-}(3,3){-}(4,4){-}(5,5)} \\
\texttt{(1,5){-}(2,4){-}(3,3){-}(4,2){-}(5,1)} \\
\texttt{So yes, (2,4) is on a main diagonal.} \\
\texttt{So all X's I have are on main diagonals. So A == "grid with X's on both main diagonals".} \\
\texttt{The user's grid is "grid with X's on both main diagonals".} \\
\texttt{So A and D are the same grid description.} \\
\texttt{The Q\&A so far has confirmed this.} \\
\texttt{So I am very confident A == D.} \\
\\ 
\texttt{The fun part is that my description was of a different grid (C/B).} \\
\texttt{But the user seems to have ignored that and is playing based on the Q\&A to compare our actual grids.} \\
\\ 
\texttt{Okay, so I'm trying to find a difference but I believe there is none.} \\
\texttt{So a good question would be one that is unambiguous and tests a specific cell.} \\
\texttt{I have asked about:} \\
\texttt{{-} (1,3) Empty} \\
\texttt{{-} (2,2) X} \\
\texttt{The user has asked about:} \\
\texttt{{-} (2,2) X} \\
\texttt{{-} Row 3: only (3,3) is X.} \\
\\ 
\texttt{What's a good, unasked{-}about cell?} \\
\texttt{(2,3) is empty.} \\
\texttt{(3,2) is empty.} \\
\texttt{(3,4) is empty.} \\
\texttt{(4,3) is empty.} \\
\texttt{(2,4) is X.} \\
\texttt{(4,2) is X.} \\
\texttt{(3,1) is empty.} \\
\texttt{(3,5) is empty.} \\
\texttt{(5,3) is empty.} \\
\texttt{(1,3) is empty. (already asked)} \\
\\ 
\texttt{User asked me: "Are there any X's in the middle row other than the center square?"} \\
\texttt{My answer: No.} \\
\texttt{This confirms my Grid A is "diagonals only".} \\
\\ 
\texttt{User answered me:} \\
\texttt{Q1: top middle empty? A: No. Correct for D.} \\
\texttt{Q2: (2,2) is X? A: Yes. Correct for D.} \\
\texttt{Q3: all diag squares are X? A: No. Correction: User answered "No" to "Are there any squares on a main diagonal that are empty?". This means ALL diag squares are X for the user. This confirms their grid is D.} \\
\\ 
\texttt{Okay, everything is 100\% consistent with me having Grid A and the user having Grid D, where A and D are identical.} \\
\texttt{So, my next question should reflect a search for a difference that I expect not to find.} \\
\\ 
\texttt{Let's ask about a square I know is X in my grid.} \\
\texttt{Maybe (4,4). This is on the (1,1){-}(5,5) diagonal.} \\
\texttt{Ok, let's ask about a square that is X in A but part of my faulty description B might be confusing.} \\
\texttt{My faulty description B had all non{-}middle{-}diagonal cells empty.} \\
\texttt{E.g. (2,4). In Grid B (plus sign), (2,4) is empty.} \\
\texttt{In Grid A/D (diagonals), (2,4) is X.} \\
\texttt{I asked about (2,2) and the user said yes.} \\
\texttt{The user already asked me about the middle row.} \\
\texttt{I haven't asked them about their middle column yet.} \\
\\ 
\texttt{Okay, new strategy: Ask about the middle column.} \\
\texttt{I know in my grid A, the middle column is ▢ ▢ X ▢ ▢.} \\
\texttt{So if the user's grid D is the same, their answer will be "no one is X except center".} \\
\texttt{If there's a difference, this will expose it.} \\
\\ 
\texttt{User's response was "ANSWER: No" to my question about an X in the middle row other than center. And then User asked me a question. The prompt *I am seeing* is asking me to answer that question, and then ask a new one. There's a small ambiguity in the provided flow.} \\
\\ 
\texttt{Let me re{-}read the last prompt again.} \\
\texttt{> ANSWER: No} \\
\texttt{> QUESTION: Are there any X's in the middle row other than the center square? Start your answer with "ANSWER:" and do not add anything else.} \\
\texttt{This is what the user sent to me.} \\
\texttt{My response should be "ANSWER: No".} \\
\texttt{Then, the next prompt I see is:} \\
\texttt{> Now ask a question in order to find out new aspects of my grid that may be different to your grid. Start with "QUESTION:" and do not add anything else.} \\
\texttt{This means my turn is to ask a question. The AI internal simulation has already happened. I am to generate the next question.} \\
\\ 
\texttt{So, I've established that user's grid is D, which is identical to my grid A.} \\
\texttt{My question needs to try to find a difference. I expect to fail.} \\
\texttt{Let's ask about (2,3). It's empty.} \\
\texttt{"QUESTION: In the second row, is the third column an X?"} \\
\texttt{This is a solid choice. It's unambiguous.} \\
\texttt{Let's try one more.} \\
\texttt{My description B was "X in the center, the four corners, and the middle of each side."} \\
\texttt{User's description D is "X's on both main diagonals."} \\
\texttt{The squares that are in my set but not theirs are (1,3), (3,1), (3,5), (5,3).} \\
\texttt{The squares that are in their set but not mine are (2,2), (2,4), (4,2), (4,4).} \\
\texttt{I already asked about the first set indirectly with "top center" {-}> user said no.} \\
\texttt{I've effectively asked about the second set with "(2,2)" {-}> user said yes.} \\
\texttt{And user asked me about their grid vs mine which seems to have confirmed they are the same.} \\
\\ 
\texttt{Okay, I'll stick with the middle column question. It's elegant.} \\
\texttt{It covers multiple cells at once.} \\
\texttt{"Are there any X's in the middle column other than the center square?"} \\
\texttt{In my grid A, the middle column is ▢ ▢ X ▢ ▢. So the answer is No.} \\
\texttt{In user's grid D, the middle column is also ▢ ▢ X ▢ ▢. So their answer will be No.} \\
\texttt{This is a good whole{-}pattern question.} \\
\\ 
\texttt{Let's consider the user's previous question to me: "Are there any X's in the middle row other than the center square?". This is a pattern question. My response was "No". So asking back about the middle column is a good parallel question. It feels fair.} \\
\\ 
\texttt{Final check:} \\
\texttt{My Grid: X ▢ ▢ ▢ X / ▢ X ▢ X ▢ / ▢ ▢ X ▢ ▢ / ▢ X ▢ X ▢ / X ▢ ▢ ▢ X} \\
\texttt{User's description: X's on both main diagonals. (Matches my grid)} \\
\texttt{User's answers:} \\
\texttt{1. (1,3) is not X. (Matches)} \\
\texttt{2. (2,2) is X. (Matches)} \\
\texttt{3. All diag squares are X. (Matches)} \\
\\ 
\texttt{Conclusion: Grids are likely identical.} \\
\texttt{My next question should test this hypothesis.} \\
\texttt{Asking about the middle column seems like a great question.} \\
\texttt{It's a high{-}information question that tests 5 cells at once (really 4, since the center is known). It's symmetrical to the last user question.} \\
\texttt{It's perfect.} \\
\texttt{Final Question Formulation:} \\
\texttt{"QUESTION: Are there any X's in the middle column other than the center square?"} \\
\texttt{This seems correct.QUESTION: Are there any X's in the middle column other than the center square?} \\
            }
        }
    }
     \\ \\

    \theutterance \stepcounter{utterance}  
    & & & \multicolumn{2}{p{0.3\linewidth}}{
        \cellcolor[rgb]{0.9,0.9,0.9}{
            \makecell[{{p{\linewidth}}}]{
                \texttt{\tiny{[GM$|$GM]}}
                \texttt{abort, first word: My, but expected ANSWER:} \\
            }
        }
    }
    & & \\ \\

    \theutterance \stepcounter{utterance}  
    & & & \multicolumn{2}{p{0.3\linewidth}}{
        \cellcolor[rgb]{0.9,0.9,0.9}{
            \makecell[{{p{\linewidth}}}]{
                \texttt{\tiny{[GM$|$GM]}}
                \texttt{Aborted} \\
            }
        }
    }
    & & \\ \\

\end{supertabular}
}

\end{document}
